El mercado internacional del vino es un entorno altamente competitivo y fragmentado. Globalmente, se producen más de 28 mil millones de litros de vino anuales repartidos en 57 países productores, sin que ninguno supere el 17\% de la participación de mercado \cite{varas2018}. En este contexto, la industria vitivinícola chilena está fuertemente orientada a la exportación, atendiendo clientes en múltiples canales de venta extranjeros \cite{varas2018}.

\section{Contexto del Negocio y Estructura del Canal\label{sec:sub1}}
Las viñas chilenas orientadas a la exportación comercializan su producción principalmente a través de grandes intermediarios, importadores internacionales o grandes cadenas de distribución minorista, más que mediante una venta directa al consumidor final \cite{bouzdine2011}. Esta dinámica concentra la negociación comercial en un grupo reducido de clientes de gran volumen. Dicha estructura se alinea con lo reportado en la literatura sobre la cadena de suministro vitivinícola nacional, la cual caracteriza a la industria exportadora por una marcada disparidad estructural entre un grupo minoritario de viñas de gran escala que concentra la mayor parte del volumen exportado, y una base atomizada de pequeños productores de uva \cite{reinecke2023}.

Esta concentración comercial simplifica las transacciones globales, pero introduce una alta complejidad a nivel operacional y de planta: cada cliente demanda habitualmente lotes en formatos estandarizados (por ejemplo, cajas de 6 o 12 botellas de 750 ml), pero bajo requerimientos estrictamente diferenciados de marca, línea comercial (económica, media o premium) y, de manera crítica para este proyecto, etiquetas específicas, a pesar de que el vino contenido en el envase sea exactamente el mismo \cite{cholette2018, varas2018}. Como consecuencia, los productos destinados a diferentes mercados geográficos no pueden ser tratados como sustitutos comerciales, a pesar de contener exactamente el mismo vino líquido \cite{cholette2009}.

\section{El Etiquetado como Problema Central\label{sec:sub2}}
Aunque el vino embotellado sea físicamente idéntico, la etiqueta debe adaptarse de manera única por al menos cuatro motivos fundamentales:

\begin{enumerate}
    \item \textbf{Idioma:} Cada mercado de destino exige su propio idioma en la etiqueta, limitando la intercambiabilidad del inventario entre mercados idiomáticos distintos \cite{bouzdine2011}.
    \item \textbf{Normativa Sanitaria y Legal:} Cada país exige advertencias distintas sobre el consumo de alcohol y menciones obligatorias de origen. En Chile, la Ley N.° 18.455 \cite{minagri1985} y su respectivo reglamento, el Decreto N.° 78 de 1986 del Ministerio de Agricultura \cite{minagri1986}, ambos fiscalizados por el Servicio Agrícola y Ganadero (SAG), exigen que los envases o etiquetas indiquen, como mínimo, la denominación, graduación alcohólica (con un mínimo de 11,5° para consumo directo), volumen y país de origen (incorporando la mención obligatoria ``Envasado en Chile'' o ``Embotellado en Chile'' para partidas de exportación) \cite{minagri1985}. Aunque el contenido conceptual sea similar entre países vecinos (por ejemplo, Chile y Argentina), la redacción legal requerida cambia, obligando a usar etiquetas físicamente diferenciadas.
    \item \textbf{Motivos Comerciales:} Los grandes compradores internacionales exigen frecuentemente el uso de marcas blancas o exclusivas, requiriendo un diseño de etiqueta diferenciado para su canal de comercialización \cite{bouzdine2011}.
    \item \textbf{Diseño de Etiqueta y Estrategia de Posicionamiento:} Las viñas adaptan los elementos estéticos (colores, tipografía y códigos visuales) a las preferencias culturales de cada mercado de destino con el fin de posicionar un mismo vino en segmentos comerciales variados (por ejemplo, estilos clásicos para mercados tradicionales versus diseños minimalistas para mercados contemporáneos) \cite{cholette2009}.
\end{enumerate}
El principal problema radica en que los pronósticos de pedidos de importación son altamente inexactos debido a la estacionalidad del negocio, los tiempos de transporte prolongados y las variaciones en las preferencias de los consumidores \cite{cholette2009, varas2018}. Planificar bajo un enfoque tradicional de fabricación para inventario (\textit{Make-to-Stock} o MTS) genera un alto riesgo de asignación errónea de producto (\textit{product misallocation}), lo que provoca simultáneamente costosas rupturas de \textit{stock} en algunos canales y excesos de inventario obsoleto en otros \cite{cholette2009, varas2018}.

\section{Comparación de Esquemas de Planificación Extremos\label{sec:sub3}}
Frente a esta multiplicidad de combinaciones vino--etiqueta, la viña puede optar por dos esquemas extremos de planificación:

\begin{itemize}
    \item \textbf{Producir contra stock (MTS)}: Fabricar y almacenar producto terminado listo para el despacho inmediato. Sin embargo, para una viña exportadora es inviable mantener \textit{stock} terminado de todas las combinaciones posibles de vino, etiqueta y mercado debido a la imprevisibilidad de los importadores y al riesgo de obsolescencia \cite{cholette2009}. Además, cambiar la etiqueta de una botella ya envasada es una operación cara y lenta, ya que implica lavar, secar y manipular botellas que ya están llenas de líquido, exponiéndolas a riesgos de rotura o mermas de calidad \cite{cholette2018}.

    \item \textbf{Producir contra pedido (MTO)}: Fabricar exclusivamente tras recibir la orden del cliente. Aunque elimina el inventario obsoleto, castiga severamente los tiempos de respuesta y satura las líneas de embotellado debido a los tiempos de preparación dependientes de la secuencia (\textit{sequence-dependent set-up times}), donde alternar la producción entre distintas familias de vinos incrementa los tiempos de lavado y merma drásticamente la capacidad de planta \cite{maccawley2022}.
\end{itemize}

\section{Planteamiento de la Oportunidad\label{sec:sub4}}
Si la viña logra desacoplar correctamente el embotellado del etiquetado, establece un punto de desacople (\textit{decoupling point}) estratégico en el inventario intermedio de botellas llenas sin etiquetar (\textit{Work-In-Process} o WIP). Este punto divide el flujo logístico en dos fases: una fase de embotellado empujada (\textit{push}) basada en pronósticos agregados de vino base, y una fase de etiquetado tirada (\textit{pull}) gatillada por los pedidos en firme de los clientes internacionales \cite{cholette2009, varas2018}. Este desacople permite reducir drásticamente el inventario terminado que hoy se mantiene ``atrapado'' en combinaciones rígidas de vino--etiqueta que no siempre calzan con la demanda real, liberando de esta manera capital de trabajo y capacidad de almacenamiento físico en la bodega \cite{cholette2009, varas2018}.

No obstante, la implementación de esta práctica introduce complejidades operacionales no menores: las líneas de producción disponen de una capacidad limitada de horas por período y se alimentan desde un número acotado de estanques intermedios, lo que limita la cantidad de variedades de vino base que pueden ser procesadas simultáneamente y obliga a secuenciar meticulosamente los ciclos para mitigar los costosos \textit{set-ups} mayores y menores \cite{maccawley2022}.

A pesar de estas restricciones, la postergación otorga una valiosa flexibilidad comercial. La bodega puede aceptar pedidos urgentes de nuevos clientes o absorber cancelaciones y modificaciones de última hora sin necesidad de reprogramar la línea de embotellado desde cero, puesto que el vino ya se encuentra envasado de forma genérica (\textit{``blanks''}) y solo resta aplicarle la etiqueta específica de destino \cite{cholette2009, varas2018}. En un contexto global donde las exigencias regulatorias y comerciales de cada mercado de destino cambian constantemente, esta flexibilidad es un factor de competitividad clave frente a viñas competidoras que producen de forma rígida contra \textit{stock} terminado y diferenciado.

Para abordar este desafío, el proyecto diseña una metodología de apoyo a la toma de decisiones basada en Investigación Operativa que permite evaluar los beneficios de postergar el etiquetado en una viña exportadora bajo incertidumbre en la demanda. La motivación, el objetivo general y los entregables comprometidos se formalizan en la Sección~\ref{sec:sec6}.